\documentclass[10pt,a4paper]{amsart}
\usepackage{amsmath,amssymb,amsthm,graphicx}
\usepackage{verbatim}
\newenvironment{reference}{\comment}{\endcomment}
%\newenvironment{reference}{}{}
\newenvironment{slogan}{\comment}{\endcomment}
\newenvironment{history}{\comment}{\endcomment}

% For commutative diagrams we use Xy-pic
\usepackage[all]{xy}

% We use 2cell for 2-commutative diagrams.
\xyoption{2cell}
% Theorem environments.
%
\theoremstyle{plain}
\newtheorem{theorem}[subsection]{Theorem}
\newtheorem{proposition}[subsection]{Proposition}
\newtheorem{lemma}[subsection]{Lemma}

\theoremstyle{definition}
\newtheorem{definition}[subsection]{Definition}
\newtheorem{example}[subsection]{Example}
\newtheorem{exercise}[subsection]{Exercise}
\newtheorem{situation}[subsection]{Situation}

\theoremstyle{remark}
\newtheorem{remark}[subsection]{Remark}
\newtheorem{remarks}[subsection]{Remarks}

\numberwithin{equation}{subsection}

% Macros
%
\def\lim{\mathop{\mathrm{lim}}\nolimits}
\def\colim{\mathop{\mathrm{colim}}\nolimits}
\def\Spec{\mathop{\mathrm{Spec}}}
\def\Hom{\mathop{\mathrm{Hom}}\nolimits}
\def\Ext{\mathop{\mathrm{Ext}}\nolimits}
\def\SheafHom{\mathop{\mathcal{H}\!\mathit{om}}\nolimits}
\def\SheafExt{\mathop{\mathcal{E}\!\mathit{xt}}\nolimits}
\def\Sch{\mathit{Sch}}
\def\Mor{\mathop{\mathrm{Mor}}\nolimits}
\def\Ob{\mathop{\mathrm{Ob}}\nolimits}
\def\Sh{\mathop{\mathit{Sh}}\nolimits}
\def\NL{\mathop{N\!L}\nolimits}
\def\CH{\mathop{\mathrm{CH}}\nolimits}
\def\proetale{{pro\text{-}\acute{e}tale}}
\def\etale{{\acute{e}tale}}
\def\QCoh{\mathit{QCoh}}
\def\Ker{\mathop{\mathrm{Ker}}}
\def\Im{\mathop{\mathrm{Im}}}
\def\Coker{\mathop{\mathrm{Coker}}}
\def\Coim{\mathop{\mathrm{Coim}}}

% Boxtimes
%
\DeclareMathSymbol{\boxtimes}{\mathbin}{AMSa}{"02}

%
% Macros for moduli stacks/spaces
%
\def\QCohstack{\mathcal{QC}\!\mathit{oh}}
\def\Cohstack{\mathcal{C}\!\mathit{oh}}
\def\Spacesstack{\mathcal{S}\!\mathit{paces}}
\def\Quotfunctor{\mathrm{Quot}}
\def\Hilbfunctor{\mathrm{Hilb}}
\def\Curvesstack{\mathcal{C}\!\mathit{urves}}
\def\Polarizedstack{\mathcal{P}\!\mathit{olarized}}
\def\Complexesstack{\mathcal{C}\!\mathit{omplexes}}
% \Pic is the operator that assigns to X its picard group, usage \Pic(X)
% \Picardstack_{X/B} denotes the Picard stack of X over B
% \Picardfunctor_{X/B} denotes the Picard functor of X over B
\def\Pic{\mathop{\mathrm{Pic}}\nolimits}
\def\Picardstack{\mathcal{P}\!\mathit{ic}}
\def\Picardfunctor{\mathrm{Pic}}
\def\Deformationcategory{\mathcal{D}\!\mathit{ef}}
\usepackage{cleveref}
\crefname{originallemma}{original lemma}{original lemmas}
\Crefname{originallemma}{Original lemma}{Original lemmas}
\makeatletter
\newlabel{decent-spaces-lemma-decent-no-specializations-map-to-same-point}{{03K5}{source}{}{sourcepointer-1}{}}
\newlabel{decent-spaces-lemma-decent-no-specializations-map-to-same-point@cref}{{[originallemma][0][]03K5}{[1][0][]source}{}{}{}}
\makeatother
\crefname{originallemma}{original lemma}{original lemmas}
\Crefname{originallemma}{Original lemma}{Original lemmas}
\makeatletter
\newlabel{decent-spaces-lemma-decent-specialization}{{03IL}{source}{}{sourcepointer-2}{}}
\newlabel{decent-spaces-lemma-decent-specialization@cref}{{[originallemma][0][]03IL}{[1][0][]source}{}{}{}}
\makeatother
\crefname{originallemma}{original lemma}{original lemmas}
\Crefname{originallemma}{Original lemma}{Original lemmas}
\makeatletter
\newlabel{duality-lemma-dimension-function-scheme}{{0AWF}{source}{}{sourcepointer-3}{}}
\newlabel{duality-lemma-dimension-function-scheme@cref}{{[originallemma][0][]0AWF}{[1][0][]source}{}{}{}}
\makeatother
\crefname{originallemma}{original lemma}{original lemmas}
\Crefname{originallemma}{Original lemma}{Original lemmas}
\makeatletter
\newlabel{dualizing-lemma-flat-unramified}{{0AWD}{source}{}{sourcepointer-4}{}}
\newlabel{dualizing-lemma-flat-unramified@cref}{{[originallemma][0][]0AWD}{[1][0][]source}{}{}{}}
\makeatother
\title{The dimension function attached to a dualizing complex on an algebraic space}
\author{The Stacks Project Authors}
\date{stacks-project:spaces-duality.tex:local-snapshot-20261012; GNU FDL 1.2 or later}
\begin{document}
\maketitle
\noindent\textit{Editorial note.} All delivered mathematical text is the author's own, copied from the pinned Stacks Project chapter file spaces-duality.tex; no statement or proof step was rewritten, reordered or removed.\par\medskip
\noindent\textit{Editorial note.} The counted proof is the author's complete argument: the comparison of the shifts of the normalized dualizing complexes through strict henselization, and the topological argument transferring the dimension function property.\par\medskip
\noindent\textit{Editorial note.} The label lemma-dimension-function-scheme names the counted statement, and the proof environment immediately following it in the source is delivered verbatim as its proof.\par\medskip
\noindent\textit{Editorial note.} The counted statement is carried by the author's own lemma environment in the chapter source.\par\medskip
\noindent\textit{Editorial note.} Cross-chapter same-name check: the label lemma-dimension-function-scheme is not carded in any other chapter. Pre-open gate: python operations/agent-authoring/tools/precard\_gate.py --source spaces-duality.tex --label lemma-dimension-function-scheme --start 485 --end 522 --local-id 28174 returned RESULT=OK.\par\medskip
\noindent\textit{Editorial note.} The statement and proof are copied from the pinned chapter file; the cited inputs are kept as external pointers with their canonical Stacks tags.\par\medskip
\noindent\textit{Allowed standard prerequisites (editorial).}\begin{itemize}
\item Behaviour of normalized dualizing complexes under unramified extension. Dualizing Complexes, the lemma describing the base change of a normalized dualizing complex along a flat unramified local homomorphism. Shows that the strict henselization of a local ring of a chart carries a normalized dualizing complex, which identifies the two candidate dimension functions.
\item The dimension function for a dualizing complex on a scheme. Duality for Schemes, the lemma attaching a dimension function on the underlying topological space to a dualizing complex on a scheme. Provides the dimension function on the chart whose composite with the map to the space is to be identified.
\item Lifting of specializations along a surjective etale morphism of decent spaces. Decent Spaces, the lemmas on specializations of decent spaces and on no nontrivial specialization mapping to the same point. Allows the topological verification that the resulting function on the space is a dimension function.
\end{itemize}\medskip
\ifdefined\hypertarget\hypertarget{sourcepointer-1}{}\fi
\noindent\textit{Source reference.} Original lemma 03K5: sources/stacks-chow/tags.txt, source lines 4665--4665. Prerequisite pointer invoked at proof line(s) 519; it names a standard fact of a sibling chapter that the counted argument uses, and every step of the counted proof built on it is retained verbatim in the delivered body.\par\medskip
\ifdefined\hypertarget\hypertarget{sourcepointer-2}{}\fi
\noindent\textit{Source reference.} Original lemma 03IL: sources/stacks-chow/tags.txt, source lines 4607--4607. Prerequisite pointer invoked at proof line(s) 518; it names a standard fact of a sibling chapter that the counted argument uses, and every step of the counted proof built on it is retained verbatim in the delivered body.\par\medskip
\ifdefined\hypertarget\hypertarget{sourcepointer-3}{}\fi
\noindent\textit{Source reference.} Original lemma 0AWF: sources/stacks-chow/tags.txt, source lines 13718--13718. Prerequisite pointer invoked at proof line(s) 512; it names a standard fact of a sibling chapter that the counted argument uses, and every step of the counted proof built on it is retained verbatim in the delivered body.\par\medskip
\ifdefined\hypertarget\hypertarget{sourcepointer-4}{}\fi
\noindent\textit{Source reference.} Original lemma 0AWD: sources/stacks-chow/tags.txt, source lines 13716--13716. Prerequisite pointer invoked at proof line(s) 506; it names a standard fact of a sibling chapter that the counted argument uses, and every step of the counted proof built on it is retained verbatim in the delivered body.\par\medskip
\begin{lemma}
\label{lemma-dimension-function-scheme}
Let $S$ be a scheme. Let $X$ be a locally Noetherian
quasi-separated algebraic space over $S$.
Let $\omega_X^\bullet$ be a dualizing complex on $X$. Then $X$ the function
$|X| \to \mathbf{Z}$ defined by
$$
x \longmapsto \delta(x)\text{ such that }
\omega_{X, \overline{x}}^\bullet[-\delta(x)]
\text{ is a normalized dualizing complex over }
\mathcal{O}_{X, \overline{x}}
$$
is a dimension function on $|X|$.
\end{lemma}

\begin{proof}
Let $U$ be a scheme and let $U \to X$ be a surjective \'etale morphism.
Let $\omega_U^\bullet$ be the dualizing complex on $U$ associated
to $\omega_X^\bullet|_U$.
If $u \in U$ maps to $x \in |X|$, then $\mathcal{O}_{X, \overline{x}}$
is the strict henselization of $\mathcal{O}_{U, u}$. By
Dualizing Complexes, Lemma \ref{dualizing-lemma-flat-unramified}
we see that if $\omega^\bullet$ is a normalized dualizing complex
for $\mathcal{O}_{U, u}$, then
$\omega^\bullet \otimes_{\mathcal{O}_{U, u}} \mathcal{O}_{X, \overline{x}}$
is a normalized dualizing complex for $\mathcal{O}_{X, \overline{x}}$.
Hence we see that the dimension function $U \to \mathbf{Z}$ of
Duality for Schemes, Lemma \ref{duality-lemma-dimension-function-scheme}
for the scheme $U$ and the complex
$\omega_U^\bullet$ is equal to the composition of $U \to |X|$ with $\delta$.
Using the specializations in $|X|$ lift to specializations in $U$
and that nontrivial specializations in $U$ map to
nontrivial specializations in $X$
(Decent Spaces, Lemmas \ref{decent-spaces-lemma-decent-specialization} and
\ref{decent-spaces-lemma-decent-no-specializations-map-to-same-point})
an easy topological argument shows that $\delta$ is a dimension function
on $|X|$.
\end{proof}

\end{document}
